\Titular*%
{miscelanea}
{Frases célebres}%
{Manuel Galán Rodríguez}%
{}

Inicio esta sección na revista Momentum achegando todas as frases que poida
entre a anterior edición da revista e a seguinte. Evidentemente, sempre se
admiten colaboracións para futuras edicións. Espero que vos guste!

\begin{multicols}{2}
\raggedcolumns

\Cita{¿Recordáis las palabras de Fernando Simón de que como mucho iba a haber 1
ó 2 casos? Pues en esta asignatura va a ser igual: va a haber como mucho 1 ó 2
suspensos.} {Diego Martínez,\\ Física computacional avanzada, 2025.}

% \Cita{Esta teoría una vez que uno la entiende, acaba reconfortado. Igual que
% uno acaba reconfortado después de haber hecho ejercicio físico.}
% {Juan Saborido,\\ Física de partículas I, 2025.}

% \Cita{Esto está explícitamente demostrado en muy pocos casos porque es
% horrible. [...] La vida es corta. Esto a los matemáticos les horroriza, pero
% después lo usan.} {Néstor,\\ Teoría cuántica de campos, 2025.}

\Cita{Igual estáis pensando que estoy haciendo una cosa extraña. Os lo va a
dejar de parecer dentro de un minuto porque la voy a complicar más.}
{Néstor Armesto,\\ Teoría cuántica de campos, 2025.}

% \Cita{Las matemáticas… las matemáticas son un lenguaje que de por sí solo no
% nos dice nada. Cuando queremos hacer poesía con ese lenguaje… lo llamamos
% Física.} {Luismi,\\ Mecánica estatística I, 2025.}

% \Cita{Que facemos? Seguimos ou xogamos unha partida de cartas?.}
% {Raúl,\\ Tecnoloxía do láser, 2025.}

\Cita{Y a mí [me da pereza] la física. Y la vida en general.}
{Víctor Pardo,\\ Electromagnetismo I, 2025.}

\Cita{El péndulo está un poco bailarín... no sabe para dónde va. Es como yo
hoy, no sé para dónde voy.}
{Jaime Álvarez,\\ Técnicas Experimentais II, 2025.}

\Cita{A potencia [solar] que chega aquí [polo norte] é moi pequena. Isto o
saben ben os pingüinos.}
{Carlos Montero,\\ Óptica I, 2025.}


\end{multicols}
